% Research contribution prepared independently for the root article.
% Needs amsmath, amssymb, amsthm.  All labels have prefix cld:.
% Suggested references: EspinosaMoll2002, DLMFhurwitz, IncomingMellinDilation,
% IncomingStieltjesClosure.  Classical spectral kernels are credited inputs.

\section{Colliding unequal dilations: a complete finite-part calculus}
\label{xid:cld:section}

The separated-grid theorem in the incoming Mellin--dilation report
\cite{IncomingMellin}
explicitly excludes $Pa\in\mathbb Z$.  At these shifts its partition of
unity no longer defines a product of the two singular distributions.
This section treats precisely that omitted case.  The construction is
the constant term of the pointwise product in a fixed ambient coordinate;
it neither assumes that a product of colliding distributions exists nor
identifies a bare spectral Laurent coefficient with that constant term.

Let $p,q$ be positive integers and put
\[
 d=\gcd(p,q),\qquad P=p/d,\qquad Q=q/d,\qquad
 A=\log P,\quad B=\log Q,\quad
 K=\log\operatorname{lcm}(p,q).
\]
For $m,n,r,k\ge0$, write $N=r+k$ and define
\begin{equation}
 J_{mn;r,k}^{p,q}
 =\operatorname{FP}_x\int_0^1
       \gamma_m^{(r)}(\{px\})\gamma_n^{(k)}(\{qx\})\,dx.
 \label{xid:cld:def:J}
\end{equation}
The derivatives are in the special-function argument.  At every member
of the union of the two singular grids, delete the right interval of
ambient length $\epsilon$ and retain the constant term as
$\epsilon\downarrow0$.  Equivalently, take separate constant terms in
independent ambient cutoff variables at the finitely many points.

Every colliding shift $a$ gives the same value as \eqref{xid:cld:def:J}.
Indeed, $Pa=h\in\mathbb Z$, and a solution of
$Qj\equiv-h\pmod P$ makes $x_0=j/p$ satisfy
$px_0,qx_0+a\in\mathbb Z$.  Translation by $x_0$ carries the two
fractional-part arguments to $px$ and $qx$ and preserves ambient lengths.
Thus no collision cases are lost by writing the common shift as zero.

\subsection{The complete local subtraction polynomial}

Use the holomorphic Stieltjes generating function
\[
 \mathscr G(z,x)=\zeta(1-z,x)+\frac1z
       =\sum_{m\ge0}\gamma_m(x)\frac{z^m}{m!},
 \qquad
 p_r(z)=\frac{(1-z)_r}{r!}
       =\prod_{j=1}^r\left(1-\frac zj\right).
\]
For $r=0$ the empty product is one.  Set
\begin{equation}
 S_{m,r}(L)=m![z^m]\,p_r(z)e^{zL}
 =\sum_{j=0}^{\min(m,r)}(-1)^j\frac{m!}{(m-j)!}
     e_j(1,1/2,\ldots,1/r)L^{m-j}.
 \label{xid:cld:S-polynomial}
\end{equation}
The Hurwitz shift identity gives the exact local decomposition
\begin{equation}
 \gamma_m^{(r)}(y)
 =(-1)^r r!y^{-r-1}S_{m,r}(\log y)
      +\gamma_m^{(r)}(1+y).
 \label{xid:cld:local-decomposition}
\end{equation}

Here is an explicit subtraction rule that covers every grid configuration.
At a point $b$ of the first grid put $t=x-b>0$ and
\[
 U_{p,m,r}(t)=(-1)^r r!p^{-r-1}t^{-r-1}
                     S_{m,r}(\log p+\log t).
\]
If $b$ is not on the second grid, subtract
\[
 U_{p,m,r}(t)\sum_{j=0}^r
       \frac{q^jt^j}{j!}\gamma_n^{(k+j)}(\{qb\}).
\]
At a point of the second grid only, use the analogous expression with
the two factors exchanged.  At a common point subtract the sum of
\begin{align}
 &U_{p,m,r}(t)U_{q,n,k}(t),\\
 &U_{p,m,r}(t)\sum_{j=0}^r
       \frac{q^jt^j}{j!}\gamma_n^{(k+j)}(1),\\
 &U_{q,n,k}(t)\sum_{j=0}^k
       \frac{p^jt^j}{j!}\gamma_m^{(r+j)}(1).
 \label{xid:cld:common-subtraction}
\end{align}
The first line has power $t^{-N-2}$ and a polynomial in $\log t$.
The other lines contain the powers $t^{-r-1},\ldots,t^{-1}$ and
$t^{-k-1},\ldots,t^{-1}$.  These are all nonintegrable terms.

Taylor's theorem and \eqref{xid:cld:local-decomposition} show that the
remaining integrand is $O(1+|\log t|^{m+n})$.  It is therefore integrable.
Adding the elementary finite-part primitives of the subtracted terms
proves existence of \eqref{xid:cld:def:J}.  For clarity, the primitive
normalization is
\[
 \operatorname{FP}\int_0^\ell t^{-1}\log^j t\,dt
      =\frac{\log^{j+1}\ell}{j+1},
\]
and, for $h\ge2$,
\[
 \operatorname{FP}\int_0^\ell t^{-h}\log^j t\,dt
 =\left.\partial_\alpha^j
       \frac{\ell^{\alpha-h+1}}{\alpha-h+1}\right|_{\alpha=0}.
\]
No finite local constant is omitted or added.

\subsection{A holomorphic master kernel for every positive derivative order}

For $M\ge1$ define
\begin{align}
 W_M(z)&=(1-z)_M\zeta(M+1-z),\\
 R(z,w)&=\frac{\Gamma(1-z-w)\Gamma(1+z)}{\Gamma(1-w)},\\
 V_M(z,w)&=R(z,w)W_M(z+w)
 =\frac{\Gamma(M+1-z-w)\Gamma(1+z)}{\Gamma(1-w)}
                  \zeta(M+1-z-w).
 \label{xid:cld:VW}
\end{align}
In particular $V_M(0,w)=W_M(w)$.  All these functions are holomorphic
in a sufficiently small polydisc about $(0,0)$.

\begin{theorem}[Colliding-grid master theorem]
\label{xid:cld:master}
For $N=r+k\ge1$, the exponential generating function of the
ambient-coordinate finite parts is
\begin{align}
 \mathscr J_{r,k}^{p,q}(z,w)
 &: =\sum_{m,n\ge0}J_{mn;r,k}^{p,q}\frac{z^mw^n}{m!n!}\notag\\
 &=Q^rP^k e^{-Bz-Aw}\left\{
 (-1)^k\frac{V_N(z,w)-e^{Kz}p_r(z)W_N(w)}{z}\right.\notag\\
 &\hspace{35mm}\left.
 +(-1)^r\frac{V_N(w,z)-e^{Kw}p_k(w)W_N(z)}{w}
                                      \right\}.
 \label{xid:cld:master-equation}
\end{align}
Both displayed divided differences have removable denominators separately.
Consequently \eqref{xid:cld:master-equation} is a finite coefficient algorithm
for every $m,n,r,k,p,q$.
\end{theorem}

\begin{proof}
First work where all derivatives of the Hurwitz factors have absolutely
convergent Fourier series.  Orthogonality leaves the frequency indices
$Q\ell$ and $-P\ell$.  The classical coincident Hurwitz product formula
\cite[Theorem 3.1]{EspinosaMoll1}, or the same Fourier calculation followed
by the Riemann functional equation, gives
\begin{align}
 \mathscr C_{r,k}^{p,q}(z,w)
 &:=\int_0^1
    \partial_y^r\zeta(1-z,\{px\})
    \partial_y^k\zeta(1-w,\{qx\})\,dx\notag\\
 &=Q^rP^ke^{-Bz-Aw}
       \left\{(-1)^k\frac{V_N(z,w)}z
                   +(-1)^r\frac{V_N(w,z)}w\right\}.
 \label{xid:cld:spectral}
\end{align}
The integral notation on the left subsequently means meromorphic
continuation.  One can verify the signs directly from
$\partial_y^r\zeta(s,y)=(-1)^r(s)_r\zeta(s+r,y)$ and
$(-1)^r(s)_r\Gamma(1-s-r)=\Gamma(1-s)$.

The local subtraction above continues the integral and identifies its
Taylor coefficients with fixed-coordinate finite parts.  We detail the
terms that distinguish these coefficients from those of
\eqref{xid:cld:spectral}.  At a first-grid point, the singular generating
factor is
\[
 (-1)^r r!p^{-r-1}p_r(z)e^{z\log p}t^{z-r-1}.
\]
Only the Taylor coefficient of degree $r$ of the other factor creates a
pole at $z=0$.  At a common point the other factor is its regular piece
at argument $1$; at a noncommon point it is its ordinary value.  These
arguments, over the first grid, form $d$ copies of
$1/P,2/P,\ldots,1$.  Since $N\ge1$, their trace is
\[
 \frac1p\sum_{b\text{ in first grid}}
      \partial_y^N\mathscr G(w,\eta_b)
   =(-1)^N P^{N-w}W_N(w),
 \qquad \eta_b=\begin{cases}\{qb\},&\{qb\}>0,\\1,&\{qb\}=0.
                         \end{cases}
\]
This is precisely the Hurwitz multiplication formula differentiated
$N$ times in the argument.

The complete $t^{z-1}$ coefficient, summed over the first grid, is
\[
 c_z(z,w)=(-1)^kQ^rP^k e^{z\log p-Aw}p_r(z)W_N(w).
\]
The meromorphic integral of $c_z t^{z-1}$ on $(0,\ell)$ is
$c_z\ell^z/z$.  The generating function of its coefficientwise
ambient finite parts is instead $c_z(\ell^z-1)/z$.
Thus the correction is the \emph{complete} quantity $-c_z(z,w)/z$,
not merely its residue at $z=0$.  The second grid contributes
\[
 -\frac{c_w(z,w)}w,
 \qquad
 c_w(z,w)=(-1)^rQ^rP^k e^{w\log q-Bz}p_k(w)W_N(z).
\]
At common points the singular--singular term is proportional to
$t^{z+w-N-2}$.  Its integral has denominator $z+w-N-1$, which is
nonzero near the origin.  Its analytic integral already equals the
generating function of its finite parts; there is no additional
diagonal contact term at $(0,0)$.

All remaining subtracted integrals converge locally uniformly on a
small spectral polydisc.  Indeed the Taylor remainders gain enough
powers of $t$, and their spectral derivatives introduce only powers
of $\log t$.  They are holomorphic by dominated differentiation.
Subtracting $c_z/z+c_w/w$ from \eqref{xid:cld:spectral} proves
\eqref{xid:cld:master-equation}, since $\log p+B=\log q+A=K$.
If one of $r,k$ is zero, replacing that spectral factor by
$\mathscr G$ adds its scalar $1/z$ or $1/w$; the integral of the other,
positive-order factor is zero by Fourier orthogonality and meromorphic
continuation.  Hence no extra scalar term occurs when $N\ge1$.
\end{proof}

The proof isolates a useful fact: common points introduce a stronger
power divergence, but the finite discrepancy between spectral and
coordinate regularization is governed by the resonant $t^{-1}$ terms.
The stronger merged term must be subtracted, yet its denominator has no
spectral pole at the parameter values being differentiated.

The comparison with rectangular spectral coefficient extraction is
completely explicit.  Define
\[
 \kappa_{m,r}(L)=m![z^{m+1}]p_r(z)e^{Lz},\qquad
 \omega_{n,N}(L)=n![w^n]e^{Lw}W_N(w),
\]
and let $\operatorname{LC}_{mn}\mathscr C$ denote
$m!n![z^mw^n]$ in the rectangular Laurent expansion of
\eqref{xid:cld:spectral}.  Then
\begin{align}
 J_{mn;r,k}^{p,q}-\operatorname{LC}_{mn}\mathscr C_{r,k}^{p,q}
 =-Q^rP^k\{&(-1)^k\kappa_{m,r}(\log p)\omega_{n,N}(-A)\notag\\
           &+(-1)^r\kappa_{n,k}(\log q)\omega_{m,N}(-B)\}.
 \label{xid:cld:spectral-contact-difference}
\end{align}
This finite formula follows by extracting the two complete local
corrections in the proof.  It specifies the difference between the two
regularizations for every index, not only for the zero-index examples.

\begin{corollary}[Finite closure in ordinary zeta derivatives]
For $N\ge1$, every $J_{mn;r,k}^{p,q}$ is a finite linear combination of
$\zeta^{(j)}(N+1)$, $0\le j\le m+n+1$.  Its coefficients lie in the
$\mathbb Q$-algebra generated by $A,B,K$, the rational generalized
harmonic numbers, and positive integer zeta values.  No unevaluated
multiple sum is required.  Writing $M=m+n$, the coefficient of the
highest derivative $\zeta^{(M+1)}(N+1)$ is
\begin{equation}
 (-1)^{M+1}N!Q^rP^k
       \left(\frac{(-1)^k}{m+1}+\frac{(-1)^r}{n+1}\right).
 \label{xid:cld:top-derivative}
\end{equation}
\end{corollary}
\begin{proof}
The finite products $p_r$ and $(1-z)_N$ have rational coefficients.
Moreover
\begin{equation}
 \log R(z,w)=\sum_{j\ge2}\frac{\zeta(j)}j
       \bigl((z+w)^j+(-1)^jz^j-w^j\bigr).
 \label{xid:cld:logR}
\end{equation}
Only total degree at most $m+n+1$ can enter a coefficient of
\eqref{xid:cld:master-equation}.  Taylor expansion of $\zeta(N+1-z-w)$
therefore uses derivatives of order at most $m+n+1$.  The vanishing
constant numerator makes division by $z$ or $w$ an exact shift of
polynomial coefficients.  The highest derivative can arise only from
the degree-$M+1$ term of $\zeta(N+1-z-w)$; the two binomial
coefficients after division by $z$ and $w$ give
\eqref{xid:cld:top-derivative}.  The contact products involve derivatives of
order at most $n$ or $m$, respectively, and cannot alter this coefficient.
\end{proof}

There is a stronger cancellation for equal Stieltjes indices after the
elementary logarithmic shifts forced by the reduced dilations.  Define
\[
 \gamma_{m,L}^{(r)}(x)=\sum_{j=0}^m\binom mj L^{m-j}
                                      \gamma_j^{(r)}(x).
\]
Its generating function is $e^{Lz}\partial_x^r\mathscr G(z,x)$.

\begin{corollary}[Odd-order diagonal collapse]
\label{xid:cld:diagonal-collapse}
If $r+k=N$ is odd, then for every $m\ge0$,
\begin{align}
 &\operatorname{FP}_x\int_0^1
   \gamma_{m,B}^{(r)}(\{px\})\gamma_{m,A}^{(k)}(\{qx\})\,dx\notag\\
 &\qquad=(-1)^kQ^rP^k
       \bigl(\kappa_{m,k}(K)-\kappa_{m,r}(K)\bigr)\omega_{m,N}(0).
 \label{xid:cld:diagonal-collapse-formula}
\end{align}
Thus only zeta derivatives of order at most $m$ remain, whereas a general
pair with these indices can require order $2m+1$.  In particular,
\begin{equation}
 \operatorname{FP}\int_0^1\gamma_m^{(r)}(x)\gamma_m^{(k)}(x)\,dx=0,
 \qquad r+k\ \text{odd},\quad m\ge\max(r,k).
 \label{xid:cld:vanishing-family}
\end{equation}
For $r=0,k=1$ the unequal-dilation identity becomes, for $m\ge1$,
\begin{align}
 &\operatorname{FP}_x\int_0^1
   \gamma_{m,B}(\{px\})\gamma_{m,A}'(\{qx\})\,dx\notag\\
 &\qquad=P(-K)^m\bigl\{\zeta^{(m)}(2)+m\zeta^{(m-1)}(2)\bigr\}.
 \label{xid:cld:diagonal-first-derivative}
\end{align}
Its $m=0$ value is $P\zeta(2)$.
\end{corollary}
\begin{proof}
Multiplying \eqref{xid:cld:master-equation} by $e^{Bz+Aw}$ removes its
external exponential.  Since $(-1)^r=-(-1)^k$, the two spectral terms
$V_N(z,w)/z$ and $V_N(w,z)/w$ have opposite diagonal coefficients.
Only the two explicit contact products remain, giving
\eqref{xid:cld:diagonal-collapse-formula}.  For $p=q=1$, one has $K=0$;
$\kappa_{m,r}(0)=0$ for $m\ge r$ because $p_r$ has degree $r$.
This proves \eqref{xid:cld:vanishing-family}.  Finally $p_1(z)-p_0(z)=-z$
gives $\kappa_{m,1}(K)-\kappa_{m,0}(K)=-K^m$, and
\[
 \omega_{m,1}(0)=(-1)^m
      \{\zeta^{(m)}(2)+m\zeta^{(m-1)}(2)\},\qquad m\ge1.
\]
These two identities prove the last formula.
\end{proof}

\subsection{Polygamma identities and a regularization diagnostic}

Since $\gamma_0=-\psi$, the zero-index specialization is particularly
short.  Put $H_0=0$ and $H_j=\sum_{h=1}^j1/h$.

\begin{corollary}[Every coincident polygamma product]
For $N=r+k\ge1$,
\begin{align}
 &\operatorname{FP}_x\int_0^1
      \psi^{(r)}(\{px\})\psi^{(k)}(\{qx\})\,dx\notag\\
 &\quad=N!Q^rP^k\left\{
 (-1)^{k+1}\bigl[\zeta'(N+1)+(H_N-H_r+K)\zeta(N+1)\bigr]\right.\notag\\
 &\hspace{37mm}\left.
 +(-1)^{r+1}\bigl[\zeta'(N+1)+(H_N-H_k+K)\zeta(N+1)\bigr]
                         \right\}.
 \label{xid:cld:polygamma}
\end{align}
In particular, for odd $N$,
\begin{equation}
 \operatorname{FP}_x\int_0^1
      \psi^{(r)}(\{px\})\psi^{(k)}(\{qx\})\,dx
 =(-1)^{k+1}N!Q^rP^k(H_k-H_r)\zeta(N+1).
 \label{xid:cld:odd-parity}
\end{equation}
These odd-order values are rational multiples of $\pi^{N+1}$ and are
unchanged by a common integer dilation of $p,q$.
\end{corollary}
\begin{proof}
At the origin $W_N(0)=N!\zeta(N+1)$ and
\[
 \partial_z V_N(0,0)
 =-N!\{\zeta'(N+1)+H_N\zeta(N+1)\},
 \qquad p_r'(0)=-H_r.
\]
Take the two removable limits in \eqref{xid:cld:master-equation}.
If $N$ is odd, $r,k$ have opposite parity, so the spectral derivatives
and all scale terms cancel.  Euler's formula for even zeta values
gives the stated rational multiples of $\pi^{N+1}$.
\end{proof}

For example,
\begin{align}
 \operatorname{FP}_x\int_0^1\psi(\{2x\})\psi'(\{3x\})\,dx
      &=\frac{\pi^2}{3},\\
 \operatorname{FP}_x\int_0^1\psi'(\{2x\})\psi'(\{3x\})\,dx
      &=12\{2\zeta'(3)+(1+2\log6)\zeta(3)\},\\
 \operatorname{FP}_x\int_0^1\psi''(\{2x\})\psi'(\{3x\})\,dx
      &=-54\zeta(4)=-\frac{3\pi^4}{5}.
 \label{xid:cld:examples}
\end{align}

The spectral continuation alone gives a different prescription.  Already
for $p=q=1,r=0,k=1$, its rectangular Laurent constant is zero, while
\[
 \operatorname{FP}\int_0^1\psi(x)\psi'(x)\,dx=\zeta(2).
\]
This can be checked without special-function product theory:
$\int_\epsilon^1\psi\psi'=\tfrac12(\psi(1)^2-\psi(\epsilon)^2)$,
and $\psi(\epsilon)=-1/\epsilon-\gamma+\zeta(2)\epsilon+O(\epsilon^2)$.
The constant is $\zeta(2)$.
More generally, substituting $z=\alpha t,w=t$ into that bare spectral
kernel gives constant
$-(\zeta(2)+\zeta'(2))(\alpha-\alpha^{-1})$.
Thus an unspecified one-parameter spectral finite part is not the
ambient-coordinate operation.  The local corrections in
\eqref{xid:cld:master-equation} remove this ambiguity before any coefficient
is extracted.

Two first-index identities illustrate the Stieltjes extension.  With
the same $P,A,K$ as above,
\begin{align}
 J_{10;0,1}^{p,q}
 &=P\left\{\frac{\zeta''(2)}2+K\zeta'(2)
            +\left(\frac{K^2}2+\log p-1\right)\zeta(2)\right\},\\
 J_{01;0,1}^{p,q}
 &=P\left\{-\frac{\zeta''(2)}2-(K+1)\zeta'(2)
            -\left(\frac{K^2}2+A\right)\zeta(2)\right\}.
 \label{xid:cld:first-jets}
\end{align}
Their sum for $p=q=1$ is $-\zeta'(2)-\zeta(2)$, in agreement with
the direct endpoint finite part of $(\gamma_1\gamma_0)'$.

\subsection{The zero-derivative Stieltjes family}

For $r=k=0$, the positive-order kernel must be supplemented by scalar
spectral-pole terms.  It is clearer to give a holomorphic formula using
the undilated coincident kernel, made explicit here:
\[
 \mathscr J_0(z,w)=\sum_{m,n\ge0}J_{mn;0,0}^{1,1}
                         \frac{z^mw^n}{m!n!},
 \qquad g(z)=\zeta(1-z)+\frac1z
          =\sum_{m\ge0}\gamma_m\frac{z^m}{m!}.
\]
The classical coincident spectral kernel is
\begin{equation}
 \mathscr C_{0,0}^{1,1}(z,w)
 =2\Gamma(z)\Gamma(w)(2\pi)^{-z-w}
       \cos\!\left(\frac{\pi(z-w)}2\right)\zeta(z+w).
 \label{xid:cld:N0-spectral}
\end{equation}
The same direct endpoint subtraction gives
\begin{equation}
 \mathscr J_0(z,w)=\mathscr C_{0,0}^{1,1}(z,w)+\frac1{zw}
                             -\frac{g(w)}z-\frac{g(z)}w.
 \label{xid:cld:J0-explicit}
\end{equation}
This grouped expression is holomorphic, as one can also see entirely
algebraically.  Put
\[
 T(z,w)=\frac{wR(z,w)+zR(w,z)}{z+w}.
\]
The numerator vanishes at $w=-z$, because $R(z,-z)=R(-z,z)=1$.
Thus $T$ is holomorphic, and $T(z,0)=T(0,w)=1$.  Applying the Riemann
functional equation to \eqref{xid:cld:N0-spectral} gives
\begin{equation}
 \mathscr J_0(z,w)=
 \frac{1-T(z,w)+(z+w)T(z,w)g(z+w)-wg(w)-zg(z)}{zw}.
 \label{xid:cld:J0-removable}
\end{equation}
The numerator vanishes on both coordinate axes, so division by $zw$
is removable.  This provides a self-contained finite coefficient rule
for the undilated zero-derivative case, without a previously tabulated
coefficient array.

Write $L_p=\log p$ and $L_q=\log q$.  Direct local subtraction, exactly
as in the master theorem, gives
\begin{align}
 \mathscr J_{0,0}^{p,q}(z,w)
 ={}&e^{-Bz-Aw}\mathscr J_0(z,w)\notag\\
 &+e^{-Aw}\frac{e^{-Bz}-e^{L_pz}}z\,g(w)
  +e^{-Bz}\frac{e^{-Aw}-e^{L_qw}}w\,g(z)\notag\\
 &+\frac{1-e^{-Bz-Aw}-e^{L_pz}+e^{L_pz-Aw}
                       -e^{L_qw}+e^{L_qw-Bz}}{zw}.
 \label{xid:cld:N0-kernel}
\end{align}
Every denominator on the right is removable; the last numerator vanishes
on both coordinate axes.  Equivalently, the full separated-grid
coefficient formula of the incoming report remains valid at collision
after replacing $I_{ij}(\theta)$ by $J_{ij;0,0}^{1,1}$ and both single
Stieltjes arguments $\theta,1-\theta$ by $1$.  This substitution is a
conclusion of local subtraction, not evaluation of a divergent integrand
at a singular shift.

For completeness, the scalar correction behind
\eqref{xid:cld:N0-kernel} is visible before extracting coefficients:
\[
 \mathscr J_{0,0}^{p,q}
 =\mathscr C_{0,0}^{p,q}+\frac1{zw}
  -\frac{e^{L_pz}}z\left(e^{-Aw}\zeta(1-w)+\frac1w\right)
  -\frac{e^{L_qw}}w\left(e^{-Bz}\zeta(1-z)+\frac1z\right).
\]
The terms in parentheses are holomorphic by cancellation of the zeta
pole.  The additional $1/(zw)$ comes from the two scalar terms used to
define $\mathscr G$.  Expanding this identity and its undilated version
proves \eqref{xid:cld:N0-kernel}.

In particular,
\begin{equation}
 \operatorname{FP}_x\int_0^1\psi(\{px\})\psi(\{qx\})\,dx
 =2\gamma_1-2\zeta(2)-2\gamma K-K^2+\log p\log q.
 \label{xid:cld:digamma}
\end{equation}
Indeed, \eqref{xid:cld:logR} gives
$T(z,w)=1+2\zeta(2)zw+O((|z|+|w|)^3)$, so
\eqref{xid:cld:J0-removable} yields $J_{00;0,0}^{1,1}=2\gamma_1-2\zeta(2)$.
This recovers the undilated coincident identity in the incoming report.
The occurrence of $K$ in the colliding unequal case is forced by the
ambient subtraction scale.

\subsection{Exact matching with the noncollision expansion}

Let
\[
 F_{mn;r,k}^{p,q}(\theta)=\FP_x\int_0^1
 \gamma_m^{(r)}(\{px\})\gamma_n^{(k)}(\{qx+a\})\,dx,
 \qquad \theta=\{Pa\}\in(0,1),
\]
denote the separated finite part. Its dependence on $a$ factors through
$\theta$ by the finite covering formula, or by Fourier orthogonality.  The next theorem resolves the collision
matching question, including every mixed derivative and Stieltjes index.
Set
\[
 C_N(z,w)=\frac{\Gamma(N+1-z-w)\Gamma(1+z)}{\Gamma(1-w)}.
\]
Unlike $V_N$, this definition is also regular for $N=0$.

\begin{theorem}[Convergent collision expansion]
\label{xid:cld:collision-expansion}
For $m,n,r,k\ge0$, define the explicit logarithmic polynomial
\begin{align}
 \mathcal P_{mn;r,k}^{p,q}(L)
 ={}&(-1)^kQ^rP^k m!n![z^mw^n]\notag\\
 &\quad e^{-Bz-Aw+wL}
       \frac{C_N(z,w)e^{zL}-e^{Kz}p_r(z)(1-w)_N}{z}.
 \label{xid:cld:collision-polynomial}
\end{align}
The quotient is holomorphic.  The polynomial has degree $m+n+1$ and
leading coefficient
\[
 \frac{(-1)^kN!Q^rP^k}{m+1}.
\]
There is an exact expansion
\begin{equation}
 F_{mn;r,k}^{p,q}(\theta)
 =\frac{\mathcal P_{mn;r,k}^{p,q}(\log\theta)}{\theta^{N+1}}
       +J_{mn;r,k}^{p,q}+\sum_{j\ge1}h_j\theta^j,
 \qquad 0<\theta<1,
 \label{xid:cld:collision-expansion-equation}
\end{equation}
whose regular power series converges locally uniformly for $|\theta|<1$.
Thus the regular constant at collision agrees exactly with the
independently defined merged-grid finite part.
\end{theorem}

\begin{proof}
The same Fourier and local-subtraction proof gives the separated kernel
by replacing $V_N(z,w)$ with
$C_N(z,w)\zeta(N+1-z-w,\theta)$ and $W_N(w)$ with
$(1-w)_N\zeta(N+1-w,\theta)$ in the first divided difference, and
by making the analogous replacements with $1-\theta$ in the second.
When $N=0$, the additional pure scalar terms containing $1/(zw)$
are independent of $\theta$ and remain unchanged; the Hurwitz terms
inside the contact numerators take the same $\theta$ or $1-\theta$
replacements. These statements follow directly from the scalar formula
preceding \eqref{xid:cld:digamma}.

Use the exact shift and Taylor identities
\begin{align*}
 \zeta(s,\theta)&=\theta^{-s}+\zeta(s,1+\theta),\\
 \zeta(s,1+\theta)&=\sum_{j\ge0}\frac{(-1)^j(s)_j}{j!}
                                      \zeta(s+j)\theta^j,\\
 \zeta(s,1-\theta)&=\sum_{j\ge0}\frac{(s)_j}{j!}
                                      \zeta(s+j)\theta^j.
\end{align*}
The first isolated power produces exactly
\eqref{xid:cld:collision-polynomial}.  The $j=0$ regular terms give the
colliding kernel proved above.  For $j\ge1$ there is no zeta pole near
the spectral origin, and the coefficient-generating function is
\begin{align}
 \mathscr H_j(z,w)
 ={}&\frac{Q^rP^ke^{-Bz-Aw}}{j!}
   \left\{(-1)^{k+j}
      \frac{V_{N+j}(z,w)-e^{Kz}p_r(z)W_{N+j}(w)}z\right.\notag\\
 &\hspace{30mm}\left.
    +(-1)^r
      \frac{V_{N+j}(w,z)-e^{Kw}p_k(w)W_{N+j}(z)}w\right\}.
 \label{xid:cld:regular-coefficients}
\end{align}
Here $h_j=m!n![z^mw^n]\mathscr H_j(z,w)$.
After division by $j!$, the gamma factors and rising factorials grow
at most polynomially in $j$ on a fixed sufficiently small spectral polydisc; their spectral
derivatives introduce powers of $\log j$.  Thus $|\theta|^j$ gives
locally uniform convergence for $|\theta|<1$, before and after any fixed
coefficient extraction.  For $N=0$, the common $j=0$ spectral pole is
removed in the complete kernel \eqref{xid:cld:N0-kernel}; it does not affect
the convergent $j\ge1$ series.

Finally the highest $L$-degree in
\eqref{xid:cld:collision-polynomial} comes from
$C_N(0,0)=N!$ and $e^{(z+w)L}/z$.  The second term in that numerator
contains $L$ only through $e^{wL}$ and cannot contribute degree
$m+n+1$.  Coefficient extraction gives the asserted leading coefficient.
\end{proof}

The lowest polynomial is already informative:
\begin{equation}
 \mathcal P_{00;r,k}^{p,q}(L)
    =(-1)^kN!Q^rP^k(L-H_N+H_r-K).
 \label{xid:cld:lowest-collision}
\end{equation}
For example,
\[
 \lim_{\theta\downarrow0}\left\{
 \operatorname{FP}_x\int_0^1
   \psi(\{2x\})\psi'(\{3x+a\})\,dx
     +\frac{2(\log\theta-1-\log6)}{\theta^2}\right\}
 =\frac{\pi^2}{3},\qquad \theta=\{2a\}.
\]
The finite term is therefore fixed by an exact convergent expansion,
including the ambient scale and every local subtraction term.

\subsection{Common dilation and an integration safeguard}

For $N\ge1$ and a further positive integer $e$, put $D=\log e$.  The reduced pair
$P,Q$ is unchanged and $K$ becomes $K+D$.  From the master kernel,
\begin{align}
 \mathscr J_{r,k}^{ep,eq}-\mathscr J_{r,k}^{p,q}
 ={}&-Q^rP^ke^{-Bz-Aw}
   \left\{(-1)^k p_r(z)W_N(w)e^{Kz}\frac{e^{Dz}-1}{z}\right.\notag\\
 &\hspace{30mm}\left.
        +(-1)^r p_k(w)W_N(z)e^{Kw}\frac{e^{Dw}-1}{w}\right\}.
 \label{xid:cld:common-scale}
\end{align}
This is a finite polynomial in $D$ at every fixed pair of Stieltjes
indices.  At zero index and odd $N$, the two terms cancel, explaining
the exact scale independence in \eqref{xid:cld:odd-parity}.  At even $N$ the
change is the negative of $D$ times the summed local $t^{-1}$ coefficient.

The incoming separated-grid theorem is valid as stated and requires no
retraction.  Its collision question can now be marked resolved by
\eqref{xid:cld:master-equation}, \eqref{xid:cld:N0-kernel}, and
\eqref{xid:cld:collision-expansion-equation}.  What must not be done during
integration is to multiply the two grid-supported contact distributions
at collision, or to keep only residues in the local spectral corrections.
The pointwise product has the unambiguous ambient finite part proved
here; its analytic representation requires the complete pole numerators.

% Verification: verify_collision.py evaluates the local-subtraction integral
% independently for nine polygamma cases; collision_validation.json records
% the floating-point residuals.  These are consistency checks, not enclosures.
